\documentclass[openany]{book}

\usepackage[margin=1in]{geometry}
\usepackage{amsmath,amsfonts,amsthm, amssymb}
\usepackage{yhmath}
\usepackage{mathrsfs}
\usepackage{mathtools}
\usepackage{xcolor}
\usepackage{graphicx}
\usepackage{comment}
\usepackage{tikz-cd}
\usepackage{quiver}
\usepackage{hyperref,cleveref}
\renewcommand{\familydefault}{ppl}
\newcommand{\curl}{\text{curl}}
\newcommand{\diverg}{\text{div}}
\newcommand{\tr}{\text{tr}}
\newcommand{\R}{\mathbb{R}}
\newcommand{\E}{\mathbb{E}}
\newcommand{\Z}{\mathbb{Z}}
\newcommand{\C}{\mathbb{C}}
\newcommand{\F}{\mathbb{F}}
\newcommand{\la}{\langle}
\newcommand{\ra}{\rangle}
\newcommand{\colim}{\text{colim}}
\DeclareMathOperator{\im}{im}
\DeclareMathOperator{\disc}{disc}
\let\oldemptyset\emptyset
\let\emptyset\varnothing
\newcommand{\tor}{\text{Tor}}
\newcommand{\id}{\text{id}}
\newcommand{\ext}{\text{Ext}}
\newcommand{\ptop}{\text{PTop}}
\newcommand{\pt}{\text{pt}}
\newcommand{\ach}{\text{Ach}}
\newcommand{\Q}{\mathbb{Q}}
\newcommand{\gal}{\text{Gal}}
\newcommand{\fraccomma}{\genfrac{}{}{0pt}{}{}{,}}
\definecolor{wikipediadarkblue}{rgb}{0.023, 0.270, 0.676}
\hypersetup{
    colorlinks,
    citecolor=black,
    filecolor=black,
    linkcolor=blue,
    urlcolor=red
}

% \hypersetup{
%     urlbordercolor={1 0 0}, % Red box for URLs
%     linkbordercolor=blue, % Blue box for internal links
%     % the values are in RGB format from 0 to 1
% }


\input{hui_r.tex}

\title{Calc 3 HW12
\\ 
\vspace{0.4cm}
\large Spring 2026}



% \author{(Please email hsun95@jh.edu if you see typos)}
% \date{\today}


\begin{document}

\maketitle

% \tableofcontents
% \newpage

\begin{prob}
    Evaluate the integral
    \[
    \int_C (2x^3 - y^3)dx + (x^3 + y^3)dy
    \]
    where $C$ is the unit circle and verify Green's Theorem in this case.
\end{prob}
\begin{proof}
    \begin{enumerate}
        \item We compute the line integral normally: we parametrize the circle $C(t)=(\cos t, \sin t), 0\leq t\leq 2\pi$.
        \begin{align*}
            \int_C(2x^3 - y^3)dx + (x^3 + y^3)dy&=\int_0^{2\pi}(2\cos^3t-\sin^3t, \cos^3t+\sin^3t)\cdot(-\sin t,\cos t)dt\\
            &=\int_0^{2\pi}
        \end{align*}
    \end{enumerate}
\end{proof}


\begin{prob}
    Verify Stokes' theorem for the upper hemisphere $z = \sqrt{1 - x^2 - y^2}, \ z \geq 0$, and the radial vector field $\vec{F}(x, y, z) = x\hat{i} + y\hat{j} + z\hat{k}$.
\end{prob}

\begin{prob}
    Show that the following vector field is conservative by constructing a potential function:\\
    $\vec{F}(x, y) = (e^x y^3 + y)\hat{i} + (3e^x y^2 + x)\hat{j}$.
\end{prob}
\begin{proof}
    A potential function can be 
    \begin{equation*}
        f=e^xy^3+xy
    \end{equation*}
\end{proof}

\begin{prob}
    \begin{enumerate}
        \item[(a)] Parameterize the unit sphere $S \subset \mathbb{R}^3$ using the spherical coordinates $(\rho, \varphi, \theta)$. Then determine the orientation of your parameterized sphere by calculating the unit normal at the north pole $(x, y, z) = (0, 0, 1)$.
        
        \item[(b)] Reparameterize $S$ so that the unit normal at the north pole points in the other direction. (c) Calculate the vector line integral of the vector field $\vec{F} = x\hat{i} + y\hat{j} + z\hat{k}$ along the equator by orienting it according to each of your two parameterizations above (really, just set the $z$-expression to $0$, thereby fixing $\varphi$, and then writing the equator curve as a function of $\theta$ alone.) Explain geometrically why this quantity is $0$ for both (indeed, all) parameterizations.
        
        \item[(c)] Calculate the flux of $\vec{F}$ through $S$, as a vector surface integral, under both parameterizations.
        
        \item[(d)] Integrate the divergence of F throughout the closed, unit ball $B$, consisting of $S$ and its interior, for each parameterization (now $\rho$ is not a constant $1$ anymore). Explain why Gauss' Theorem is not violated in both cases, even though your two parameterizations in part (d) will yield different answers.
    \end{enumerate}
\end{prob}
\begin{proof}
    
\end{proof}

% \begin{enumerate}
%     \item Evaluate the integral
%     \[
%     \int_C (2x^3 - y^3)dx + (x^3 + y^3)dy
%     \]
%     where $C$ is the unit circle and verify Green's Theorem in this case.

%     \item Verify Stokes' theorem for the upper hemisphere $z = \sqrt{1 - x^2 - y^2}, \ z \geq 0$, and the radial vector field $\vec{F}(x, y, z) = x\hat{i} + y\hat{j} + z\hat{k}$.

%     \item Show that the following vector field is conservative by constructing a potential function:\\
%     $\vec{F}(x, y) = (e^x y^3 + y)\hat{i} + (3e^x y^2 + x)\hat{j}$.

%     \item 
%     \begin{enumerate}
%         \item[(a)] Parameterize the unit sphere $S \subset \mathbb{R}^3$ using the spherical coordinates $(\rho, \varphi, \theta)$. Then determine the orientation of your parameterized sphere by calculating the unit normal at the north pole $(x, y, z) = (0, 0, 1)$.
        
%         \item[(b)] Reparameterize $S$ so that the unit normal at the north pole points in the other direction. (c) Calculate the vector line integral of the vector field $\vec{F} = x\hat{i} + y\hat{j} + z\hat{k}$ along the equator by orienting it according to each of your two parameterizations above (really, just set the $z$-expression to $0$, thereby fixing $\varphi$, and then writing the equator curve as a function of $\theta$ alone.) Explain geometrically why this quantity is $0$ for both (indeed, all) parameterizations.
        
%         \item[(c)] Calculate the flux of $\vec{F}$ through $S$, as a vector surface integral, under both parameterizations.
        
%         \item[(d)] Integrate the divergence of F throughout the closed, unit ball $B$, consisting of $S$ and its interior, for each parameterization (now $\rho$ is not a constant $1$ anymore). Explain why Gauss' Theorem is not violated in both cases, even though your two parameterizations in part (d) will yield different answers.
%     \end{enumerate}
% \end{enumerate}


\end{document}